% =============================================================================
% Comandos do template
% Usuários normalmente não precisam alterar este arquivo.
% =============================================================================
\makeatletter

% ---------------------------------------------------------------------------
% Metadados (preenchidos em config/metadata.tex)
% ---------------------------------------------------------------------------
\newcommand{\tcc@titulo}{}
\newcommand{\tcc@tituloingles}{}
\newcommand{\tcc@curso}{}
\newcommand{\tcc@instituicao}{}
\newcommand{\tcc@cidade}{}
\newcommand{\tcc@ano}{\the\year}
\newcommand{\tcc@palavraschave}{}
\newcommand{\tcc@keywords}{}
\newcommand{\tcc@autores}{}
\newcommand{\tcc@nomesautores}{}
\newcommand{\tcc@orientador}{}

\newcommand{\titulo}[1]{\renewcommand{\tcc@titulo}{#1}}
\newcommand{\tituloingles}[1]{\renewcommand{\tcc@tituloingles}{#1}}
\newcommand{\curso}[1]{\renewcommand{\tcc@curso}{#1}}
\newcommand{\instituicao}[1]{\renewcommand{\tcc@instituicao}{#1}}
\newcommand{\cidade}[1]{\renewcommand{\tcc@cidade}{#1}}
\newcommand{\ano}[1]{\renewcommand{\tcc@ano}{#1}}
\newcommand{\palavraschave}[1]{\renewcommand{\tcc@palavraschave}{#1}}
\newcommand{\keywords}[1]{\renewcommand{\tcc@keywords}{#1}}

% \autor[ORCID]{Nome}{Formação, vínculo}{e-mail}
% (ORCID e e-mail só aparecem quando informados: nada é inventado)
\newcommand{\autor}[4][]{%
  \appto\tcc@autores{\tcc@blocoautor{#1}{#2}{#3}{#4}}%
  \ifdefempty{\tcc@nomesautores}
    {\appto\tcc@nomesautores{#2}}
    {\appto\tcc@nomesautores{, #2}}%
}

% \orientador[ORCID]{Nome}{Titulação, vínculo}{e-mail}
\newcommand{\orientador}[4][]{%
  \renewcommand{\tcc@orientador}{\tcc@blocoautor{#1}{#2}{#3}{#4}}%
}

% Acesso aos metadados dentro do texto
\newcommand{\imprimirtitulo}{\tcc@titulo}
\newcommand{\imprimirtituloingles}{\tcc@tituloingles}
\newcommand{\imprimircurso}{\tcc@curso}
\newcommand{\imprimirinstituicao}{\tcc@instituicao}
\newcommand{\imprimircidade}{\tcc@cidade}
\newcommand{\imprimirano}{\tcc@ano}

% ---------------------------------------------------------------------------
% Primeira página: títulos (centralizados, 14 pt) e autores (à esquerda).
% Layout dos artigos FAESA aprovados (.ref/): autores à direita, nome em
% maiúsculas, formação/vínculo e e-mail sublinhado abaixo, sem nota de rodapé.
% ---------------------------------------------------------------------------
\newcommand{\tcc@blocoautor}[4]{%
  \begingroup
    \raggedleft
    \setlength{\parskip}{0pt}%
    \setstretch{1}%
    \MakeUppercase{#2}\par
    #3\par
    \ifstrempty{#1}{}{Orcid: #1\par}%
    \ifstrempty{#4}{}{\mbox{\underline{\href{mailto:#4}{#4}}}\par}%
  \endgroup
  \vspace{6pt}%
}

\newcommand{\imprimircabecalho}{%
  \begingroup
    \setlength{\parskip}{0pt}%
    \setstretch{1}%
    \centering\fontsize{14}{16.1}\selectfont\bfseries
    \MakeUppercase{\tcc@titulo}\par
    \vspace{14pt}%
    \ifdefempty{\tcc@tituloingles}{}{\itshape\MakeUppercase{\tcc@tituloingles}\par}%
  \endgroup
  \vspace{18pt}%
  \tcc@autores
  \tcc@orientador
}

% ---------------------------------------------------------------------------
% Resumo e Abstract (palavras-chave vêm de config/metadata.tex)
% Títulos sem numeração: centralizados (Guia p. 10)
% ---------------------------------------------------------------------------
\titleformat{name=\section,numberless}[block]{\normalfont\bfseries\centering}{}{0pt}{\MakeUppercase}

\newcommand{\tcc@palavras}[2]{%
  \ifdefempty{#2}{}{\par\textbf{#1:} #2\par}%
}

\newenvironment{resumo}
  {\section*{Resumo}}
  {\tcc@palavras{Palavras-chave}{\tcc@palavraschave}}

\renewenvironment{abstract}
  {\section*{Abstract}\itshape}
  {\par\textbf{Keywords:} \tcc@keywords\par}

% ---------------------------------------------------------------------------
% Figuras
% \figura[largura]{arquivo}{Título}{Fonte}{rótulo}
% ---------------------------------------------------------------------------
\newcommand{\fonte}[1]{%
  \par\vspace{\espacofonte}%
  {\raggedright\fontedez Fonte: #1\par}%
}

\newcommand{\figura}[5][\largurapadraofigura]{%
  \begin{figure}
    \caption{#3}\label{#5}
    \centering
    \includegraphics[width=#1,height=\alturamaximafigura,keepaspectratio]{#2}%
    \par
    \fonte{#4}
  \end{figure}%
}

% ---------------------------------------------------------------------------
% Citação direta longa (mais de três linhas): Arial 10, recuo de 4 cm
% ---------------------------------------------------------------------------
\newenvironment{citacao}
  {\par\begingroup
   \setlength{\leftskip}{\recuocitacao}%
   \setlength{\parskip}{0pt}%
   \fontedez}
  {\par\endgroup}

% ---------------------------------------------------------------------------
% Lista de alíneas: a), b), c)
% ---------------------------------------------------------------------------
\newlist{alineas}{enumerate}{1}
\setlist[alineas]{label=\alph*)}

% ---------------------------------------------------------------------------
% Marcadores de pendência (aparecem em vermelho no PDF; remova antes da entrega)
% \preencher{o que falta}   informação a ser fornecida pelos autores
% \verificar{o que conferir} referência ou dado a conferir no original
% ---------------------------------------------------------------------------
\newcommand{\preencher}[1]{\textcolor{red}{[A PREENCHER: #1]}}
\newcommand{\verificar}[1]{\textcolor{red}{[VERIFICAR: #1]}}
\newcommand{\definir}[1]{\textcolor{red}{[DEFINIR: #1]}}

% Algoritmos em pseudocódigo: flutuante próprio, numerado à parte dos quadros
\DeclareCaptionType[fileext=loa,placement=htbp]{algoritmo}[Algoritmo][Lista de algoritmos]

% ---------------------------------------------------------------------------
% Referências bibliográficas
% ---------------------------------------------------------------------------
\bibliographystyle{abntex2-alf}

\renewcommand{\refname}{Referências}
\AfterTheBibliography{%
  \fontedozesimples
  \raggedright
  \setlength{\parskip}{0pt}%
  \setlength{\parsep}{0pt}%
  \setlength{\itemsep}{\espacoentrereferencias}%
}

\newcommand{\imprimirreferencias}{%
  \bibliography{references/references}%
}

% Apêndices (após as referências): título centralizado, sem numeração progressiva
\newcounter{apendice}
\renewcommand{\theapendice}{\Alph{apendice}}
\newcommand{\apendice}[2]{\FloatBarrier\refstepcounter{apendice}\section*{Apêndice \theapendice\ - #1}\label{#2}}

% ---------------------------------------------------------------------------
% Metadados do PDF
% ---------------------------------------------------------------------------
\AtBeginDocument{%
  \hypersetup{
    pdftitle={\tcc@titulo},
    pdfauthor={\tcc@nomesautores},
    pdfsubject={\tcc@curso\ -- \tcc@instituicao},
    pdfkeywords={\tcc@palavraschave},
  }%
}

\makeatother
